\documentclass[a4paper, 12pt]{article}

% --- PAQUETES DE IDIOMA Y CODIFICACIÓN ---
\usepackage[utf8]{inputenc}
\usepackage[T1]{fontenc}
\usepackage[spanish, es-tabla]{babel}

% --- FORMATO APA (GEOMETRÍA Y FUENTE) ---
\usepackage[margin=2.54cm]{geometry} % Márgenes APA
\usepackage{mathptmx}                % Fuente Times New Roman
\usepackage{setspace}                % Para interlineado
\usepackage{hanging}                 % Para sangría francesa
\usepackage{listings}

% --- PAQUETES ADICIONALES ---
\usepackage{graphicx}
\usepackage{float}      % Permite usar [H]
\usepackage{amsmath}
\usepackage{amsfonts}
\usepackage{amssymb}
\usepackage{xurl}       % Corta las URLs correctamente (ya incluye url)
\usepackage{eso-pic}    % Para imagen de fondo
\usepackage{fancyhdr}   % Para encabezados personalizados
\usepackage{hyperref}   % Mejor cargarlo al final

% --- CONFIGURACIÓN DE PIES DE FOTO ---
\usepackage[
   labelfont=bf,
   font=it,
   labelsep=newline,
   justification=raggedright,
   singlelinecheck=false
]{caption}

% --- CONFIGURACIÓN DE LOS BLOQUES DE CÓDIGO ---
\lstset{
  basicstyle=\ttfamily\footnotesize,
  keywordstyle=\bfseries,
  frame=single,
  breaklines=true,        % Corta las líneas largas en vez de salirse del margen
  numbers=left,
  numberstyle=\tiny,
  showstringspaces=false,
  tabsize=2,
  captionpos=b,
  % Sin esto, las tildes y la ñ dentro del código dan error
  literate={á}{{\'a}}1 {é}{{\'e}}1 {í}{{\'i}}1 {ó}{{\'o}}1 {ú}{{\'u}}1
           {Á}{{\'A}}1 {É}{{\'E}}1 {Í}{{\'I}}1 {Ó}{{\'O}}1 {Ú}{{\'U}}1
           {ñ}{{\~n}}1 {Ñ}{{\~N}}1 {¿}{{?`}}1 {¡}{{!`}}1 {·}{{$\cdot$}}1
}
\renewcommand{\lstlistingname}{Código}

% TypeScript no viene incluido en listings
\lstdefinelanguage{TypeScript}{
  keywords={import, from, export, class, interface, type, const, private, public,
    readonly, async, await, return, if, else, throw, true, false, null, undefined,
    constructor, implements, this, string, number, Promise, new},
  sensitive=true,
  comment=[l]{//},
  morecomment=[s]{/*}{*/},
  morestring=[b]', morestring=[b]"
}

% --- CONFIGURACIÓN DE PÁRRAFOS ---
\setlength{\parindent}{0.5cm}
\setlength{\headheight}{45pt}

% --- CONFIGURACIÓN DE IMÁGENES DE FONDO Y ENCABEZADO ---
\newcommand\BackgroundPic{
    \put(0,0){
        \parbox[b][\paperheight]{\paperwidth}{
            \vfill
            \centering
            \includegraphics[width=0.7\paperwidth,keepaspectratio]{logonebrija.png}
            \vfill
        }
    }
}

\pagestyle{fancy}
\fancyhf{}
\fancyhead[R]{\thepage}
\fancyhead[L]{\includegraphics[height=1.1cm]{logoencabezado.png}}
\renewcommand{\headrulewidth}{0pt}

%---------------------------------------------------------------------------
\begin{document}

\AddToShipoutPicture*{\BackgroundPic}

% --- PORTADA ---
\begin{titlepage}
    \centering
    \vspace*{2cm}

    {\Huge \textbf{Supabase}}\\[0.5cm]
    {\Large \textbf{Desarrollo de interfaces}}\\[1.5cm]

    {\large Noelia Cantador Bosqued}\\[0.5cm]
    {\large Desarrollo de Aplicaciones Multiplataforma}\\
    {\large Institutos Nebrija}\\[2cm]

    \begin{minipage}{0.85\textwidth}
        \centering
        \textit{Gestión completa CRUD con roles (Ionic, Angular Standalone y Supabase)}
    \end{minipage}

    \vfill
    {\large \today}
    \vspace*{2cm}
\end{titlepage}

% --- ÍNDICE ---
\newpage
\tableofcontents
\newpage

% ===========================================================================
\section{Introducción}
% ===========================================================================

En este laboratorio se desarrolla una aplicación de gestión de clientes con
\textbf{Ionic} y \textbf{Angular Standalone}, conectada a un \textit{backend}
\textbf{Supabase}. La aplicación permite:

\begin{itemize}
  \item Consultar el listado de clientes y la galería de imágenes sin iniciar
        sesión.
  \item Crear, editar y borrar clientes (\textbf{CRUD}). Al intentarlo, la
        aplicación pide iniciar sesión y solo lo permite a los administradores.
  \item Ver como \textit{cards} las imágenes guardadas en un \textit{bucket} de
        Supabase Storage.
\end{itemize}

La conexión se realiza con la librería oficial \texttt{@supabase/supabase-js},
que crea un cliente a partir de la URL del proyecto y la clave pública
(\textit{publishable key}). Supabase muestra los pasos para Ionic Angular en su
asistente de conexión (figura \ref{fig:connect}).

\begin{figure}[H]
    \centering
    \includegraphics[width=0.75\linewidth]{image.png}
    \caption{Asistente de conexión de Supabase para Ionic Angular.}
    \label{fig:connect}
\end{figure}

El código completo está publicado en GitHub:
\url{https://github.com/NoeCB/supabase}.

\subsection{Resultados de aprendizaje}

\begin{itemize}
  \item \textbf{RA 1 -- CE 1.c}: enlace entre el modelo y la vista mediante
        formularios reactivos y \textit{signals}.
  \item \textbf{RA 3 -- CE 3.d}: gestión de eventos de autenticación y
        modificación de la interfaz según el rol del usuario.
  \item \textbf{RA 4 -- CE 4.a}: diseño responsivo con \texttt{ion-list} e
        \texttt{ion-grid}.
\end{itemize}

\subsection{Base teórica}

La aplicación aplica un \textbf{control de acceso basado en roles} (RBAC). En
una arquitectura \textit{serverless} con Supabase, la seguridad no está en la
aplicación Angular, sino en las \textbf{políticas de seguridad a nivel de fila}
(\textit{Row Level Security}, RLS) de PostgreSQL. La interfaz tiene un doble
papel:

\begin{enumerate}
  \item \textbf{Formularios reactivos} (\texttt{ReactiveFormsModule}):
        validan los datos en el cliente antes de enviarlos.
  \item \textbf{Renderizado condicional} (\texttt{@if}): adapta la interfaz al
        estado de la sesión, por ejemplo mostrando el botón de iniciar o cerrar
        sesión y el rol del usuario.
\end{enumerate}

En la aplicación hay tres niveles de acceso:

\begin{table}[H]
\centering
\begin{tabular}{|l|l|l|}
\hline
\textbf{Usuario} & \textbf{Perfil en \texttt{profiles}} & \textbf{Permisos} \\ \hline
Visitante (sin sesión) & --             & Ver clientes y galería \\ \hline
noeliac@gmail.com      & sin perfil     & Ver, pero no modificar clientes \\ \hline
juan@gmail.com         & \texttt{admin} & Ver, crear, editar y borrar clientes \\ \hline
\end{tabular}
\caption{Niveles de acceso de la aplicación.}
\label{tab:roles}
\end{table}

% ===========================================================================
\section{Paso 1: Instalación de la librería}
% ===========================================================================

El proyecto parte de una aplicación Ionic con Angular Standalone (plantilla de
pestañas). Desde la terminal integrada, en la raíz del proyecto, se instala el
SDK oficial de Supabase:

\begin{lstlisting}[language=bash, caption={Instalación del cliente de Supabase}]
npm install @supabase/supabase-js
\end{lstlisting}

Esto añade la dependencia al archivo \texttt{package.json} del proyecto.

\begin{figure}[H]
    \centering
    \includegraphics[width=0.75\linewidth]{image1.png}
    \caption{Resultado de la instalación de la dependencia.}
    \label{fig:npm-install}
\end{figure}

% ===========================================================================
\section{Paso 2: Configuración de las credenciales (environments)}
% ===========================================================================

La URL y la clave del proyecto se guardan en los archivos de entorno de
Angular, separadas de la lógica de la aplicación. Se usa la
\textbf{publishable key}, que es pública y está pensada para el cliente: lo que
protege los datos son las políticas RLS. La clave \texttt{service\_role} nunca
debe aparecer en el código, porque se salta todas las políticas.

\subsection*{src/environments/environment.ts}

\begin{lstlisting}[language=TypeScript, caption={Credenciales en modo desarrollo}]
export const environment = {
  production: false,
  supabaseUrl: 'https://rueljvohvzfdmtfjwazl.supabase.co',
  supabaseKey: 'sb_publishable_AhKlJL6ZZwxTRmuCQVOMSg_j2oCc86e',
};
\end{lstlisting}

\subsection*{src/environments/environment.prod.ts}

\begin{lstlisting}[language=TypeScript, caption={Credenciales en modo producción}]
export const environment = {
  production: true,
  supabaseUrl: 'https://rueljvohvzfdmtfjwazl.supabase.co',
  supabaseKey: 'sb_publishable_AhKlJL6ZZwxTRmuCQVOMSg_j2oCc86e',
};
\end{lstlisting}

La librería añade por su cuenta la ruta de la API REST
(\texttt{/rest/v1/}), así que en la URL solo se pone el dominio del proyecto.

% ===========================================================================
\section{Paso 3: Configuración del backend en Supabase (SQL)}
% ===========================================================================

Todo el SQL se ejecuta en el \textbf{SQL Editor} de Supabase
(icono \texttt{>\_} de la barra lateral). Cuando hay varias sentencias, hay
que ejecutar el bloque completo; si solo se selecciona una parte, el editor
ejecuta únicamente esa parte.

\subsection{Tablas, roles y políticas RLS}

El script del laboratorio crea el tipo de rol, las tablas \texttt{profiles} y
\texttt{clients} y las políticas de seguridad:

\begin{lstlisting}[language=SQL, caption={Script de creación de la base de datos}]
-- 1. Enumeración de roles
CREATE TYPE public.user_role AS ENUM ('admin', 'user');

-- 2. Tabla de perfiles (vinculada a autenticación)
CREATE TABLE public.profiles (
  id UUID REFERENCES auth.users(id) ON DELETE CASCADE PRIMARY KEY,
  role public.user_role DEFAULT 'user' NOT NULL
);

-- 3. Tabla de clientes
CREATE TABLE public.clients (
  id UUID DEFAULT gen_random_uuid() PRIMARY KEY,
  nombre TEXT NOT NULL,
  empresa TEXT NOT NULL,
  email TEXT NOT NULL UNIQUE,
  created_at TIMESTAMPTZ DEFAULT NOW()
);

-- 4. Habilitar RLS
ALTER TABLE public.clients ENABLE ROW LEVEL SECURITY;
ALTER TABLE public.profiles ENABLE ROW LEVEL SECURITY;

-- 5. Cada usuario solo puede leer su propio perfil
CREATE POLICY "Leer perfil propio" ON public.profiles
  FOR SELECT USING (auth.uid() = id);

-- 6. Lectura de clientes: usuarios autenticados
CREATE POLICY "Lectura para autenticados" ON public.clients
  FOR SELECT TO authenticated USING (true);

-- 7. Escritura (insert, update, delete): solo administradores
CREATE POLICY "Escritura solo para admins" ON public.clients
  FOR ALL TO authenticated
  USING (EXISTS (SELECT 1 FROM public.profiles
                 WHERE id = auth.uid() AND role = 'admin'))
  WITH CHECK (EXISTS (SELECT 1 FROM public.profiles
                      WHERE id = auth.uid() AND role = 'admin'));
\end{lstlisting}

El resultado es el esquema de la figura \ref{fig:esquema}: \texttt{profiles}
está relacionada con \texttt{auth.users} por su \texttt{id}, y
\texttt{clients} guarda los datos de los clientes.

\begin{figure}[H]
    \centering
    \includegraphics[width=0.75\linewidth]{image3.png}
    \caption{Base de datos de clients/profiles}
    \label{fig:esquema}
\end{figure}

\subsection{Lectura pública de los clientes}

Con el script original, solo los usuarios con sesión pueden leer los clientes
(\texttt{TO authenticated}). Como la aplicación debe mostrar el listado a
cualquier visitante, se sustituye la política de lectura por otra que también
incluye el rol \texttt{anon} (usuario sin sesión):

\begin{lstlisting}[language=SQL, caption={Política de lectura pública de clients}]
DROP POLICY IF EXISTS "Lectura para autenticados" ON public.clients;
CREATE POLICY "Lectura publica" ON public.clients
  FOR SELECT TO anon, authenticated
  USING (true);
\end{lstlisting}

La política \textit{``Escritura solo para admins''} no se modifica: crear,
editar y borrar sigue reservado a los usuarios con rol \texttt{admin}.

\subsection{Creación de usuarios y asignación de roles}

Los usuarios se crean en \textit{Authentication > Users > Add user >
Create new user}, marcando \textit{Auto Confirm User} para no tener que
confirmar el correo. Se crean dos usuarios:

\begin{itemize}
  \item \texttt{juan@gmail.com}, que será el administrador.
  \item \texttt{noeliac@gmail.com}, que podrá iniciar sesión pero no modificar
        datos.
\end{itemize}

El script del laboratorio no crea los perfiles automáticamente, así que el rol
se asigna a mano. Juan recibe el rol \texttt{admin} y Noelia se queda sin
perfil, por lo que la política de escritura la bloquea:

\begin{lstlisting}[language=SQL, caption={Asignación de roles}]
-- Juan es administrador
INSERT INTO public.profiles (id, role)
SELECT id, 'admin' FROM auth.users WHERE email = 'juan@gmail.com'
ON CONFLICT (id) DO UPDATE SET role = 'admin';

-- Noelia sin perfil = sin permisos
DELETE FROM public.profiles
WHERE id = (SELECT id FROM auth.users WHERE email = 'noeliac@gmail.com');
\end{lstlisting}

Para comprobar el rol de cada usuario se cruza \texttt{auth.users} con
\texttt{profiles}. Juan debe aparecer como \texttt{admin} y Noelia con
\texttt{NULL}:

\begin{lstlisting}[language=SQL, caption={Comprobación de los roles}]
SELECT u.email, p.role
FROM auth.users u
LEFT JOIN public.profiles p ON p.id = u.id;
\end{lstlisting}

Un usuario con rol \texttt{user} tendría los mismos permisos que Noelia: puede
leer, pero no escribir, porque la política de escritura exige el rol
\texttt{admin}.

\subsection{Datos de clientes}

Los clientes se pueden crear desde la propia aplicación (sesión de Juan) o
desde el SQL Editor. El \texttt{email} es único, así que no se puede repetir:

\begin{lstlisting}[language=SQL, caption={Inserción de clientes de ejemplo}]
INSERT INTO public.clients (nombre, empresa, email) VALUES
  ('María López',   'Tech Solutions SL', 'maria.lopez@ejemplo.com'),
  ('Carlos Ruiz',   'Innova Digital',    'carlos.ruiz@ejemplo.com'),
  ('Laura Gómez',   'Grupo Norte SA',    'laura.gomez@ejemplo.com');
\end{lstlisting}

\subsection{Storage: bucket de imágenes}

En \textit{Storage > Files} se crea el \textit{bucket} privado
\texttt{galeria} y se suben las imágenes con \textit{Upload files}. Aunque el
\textit{bucket} tenga archivos, la aplicación no los ve si no hay una
política que lo permita: Supabase devuelve una lista vacía, no un error. Por
eso se añade una política de lectura pública sobre \texttt{storage.objects},
limitada a este \textit{bucket}:

\begin{lstlisting}[language=SQL, caption={Política de lectura del bucket galeria}]
CREATE POLICY "Ver galeria publica" ON storage.objects
  FOR SELECT TO anon, authenticated
  USING (bucket_id = 'galeria');
\end{lstlisting}

Solo se da permiso de lectura (\texttt{SELECT}): nadie puede subir ni borrar
imágenes desde la aplicación.

Las políticas activas de Storage se pueden consultar con:

\begin{lstlisting}[language=SQL, caption={Consulta de las políticas de Storage}]
SELECT policyname, cmd, roles, qual
FROM pg_policies
WHERE schemaname = 'storage' AND tablename = 'objects';
\end{lstlisting}

% ===========================================================================
\section{Paso 4: Modelos de datos}
% ===========================================================================

Las interfaces de TypeScript reflejan la forma de los datos de Supabase, de
modo que el compilador detecta si se usa un campo que no existe.

\begin{lstlisting}[language=TypeScript, caption={src/app/models/client.model.ts}]
export interface Client {
  id?: string; // UUID
  nombre: string;
  empresa: string;
  email: string;
  created_at?: string;
}

export type UserRole = 'admin' | 'user';
\end{lstlisting}

\begin{lstlisting}[language=TypeScript, caption={src/app/models/storage-image.model.ts}]
export interface StorageImage {
  name: string;
  url: string;   // URL firmada (el bucket es privado)
  size?: number; // bytes
  createdAt?: string;
}
\end{lstlisting}

% ===========================================================================
\section{Paso 5: Servicio de conexión (supabase.service.ts)}
% ===========================================================================

El servicio \texttt{SupabaseService} (\texttt{providedIn: 'root'}, una única
instancia para toda la aplicación) crea el cliente de Supabase y centraliza la
autenticación, los roles, el CRUD y el acceso a Storage.

\subsection{Conexión, sesión y roles}

El usuario y su rol se guardan en \textit{signals}. Cuando cambian, la
interfaz se actualiza sola. El servicio recupera la sesión guardada al
arrancar y escucha los cambios de sesión con \texttt{onAuthStateChange}.

\begin{lstlisting}[language=TypeScript, caption={Conexión y estado de la sesión}]
export class SupabaseService {
  private readonly supabase: SupabaseClient;

  // Estado global reactivo (Signals)
  readonly currentUser = signal<User | null>(null);
  readonly userRole = signal<UserRole | null>(null);
  readonly isAuthenticated = computed(() => this.currentUser() !== null);
  readonly isAdmin = computed(() => this.userRole() === 'admin');

  // Se resuelve cuando la sesión guardada se ha recuperado
  readonly ready: Promise<void>;

  constructor() {
    this.supabase = createClient(environment.supabaseUrl, environment.supabaseKey);

    this.ready = this.supabase.auth
      .getSession()
      .then(({ data }) => this.applySession(data.session));

    // Mantiene el estado sincronizado (login, logout, refresco de token)
    this.supabase.auth.onAuthStateChange((_event, session) => {
      setTimeout(() => this.applySession(session), 0);
    });
  }

  private async loadRole(userId: string): Promise<void> {
    const { data, error } = await this.supabase
      .from('profiles').select('role').eq('id', userId).maybeSingle();
    if (error) console.error('Error cargando el rol:', error.message);
    // Sin perfil en la tabla profiles = sin permisos (null)
    this.userRole.set((data?.role as UserRole) ?? null);
  }

  async login(email: string, password: string): Promise<void> {
    const { data, error } = await this.supabase.auth
      .signInWithPassword({ email, password });
    if (error) throw error;
    await this.applySession(data.session);
  }
}
\end{lstlisting}

\subsection{CRUD de clientes}

\begin{lstlisting}[language=TypeScript, caption={Operaciones CRUD sobre la tabla clients}]
async getClients(): Promise<Client[]> {
  const { data, error } = await this.supabase
    .from('clients').select('*')
    .order('created_at', { ascending: false });
  if (error) throw error;
  return data ?? [];
}

async addClient(client: Client): Promise<void> {
  const { error } = await this.supabase.from('clients').insert({
    nombre: client.nombre, empresa: client.empresa, email: client.email,
  });
  if (error) throw error;
}

async updateClient(id: string, client: Partial<Client>): Promise<void> {
  const { data, error } = await this.supabase.from('clients')
    .update({ nombre: client.nombre, empresa: client.empresa, email: client.email })
    .eq('id', id)
    .select();
  if (error) throw error;
  // Si RLS bloquea la operación, Supabase no da error: no modifica ninguna fila
  if (!data?.length) throw new Error('No tienes permisos para editar este cliente');
}

async deleteClient(id: string): Promise<void> {
  const { data, error } = await this.supabase
    .from('clients').delete().eq('id', id).select();
  if (error) throw error;
  if (!data?.length) throw new Error('No tienes permisos para borrar este cliente');
}
\end{lstlisting}

\subsection{Imágenes del bucket}

Como el \textit{bucket} es privado, no sirven las URLs públicas. El servicio
lista los archivos, descarta lo que no sea una imagen y genera URLs firmadas
que caducan a la hora:

\begin{lstlisting}[language=TypeScript, caption={Lectura de las imágenes de Storage}]
async getGalleryImages(): Promise<StorageImage[]> {
  const bucket = this.supabase.storage.from('galeria');

  // 1. Listar los archivos de la raíz del bucket
  const { data: files, error } = await bucket.list('', {
    limit: 100,
    sortBy: { column: 'name', order: 'asc' },
  });
  if (error) throw error;

  // Solo imágenes (se descartan carpetas y marcadores)
  const images = (files ?? []).filter((f) => f.id && IMAGE_EXTENSIONS.test(f.name));
  if (!images.length) return [];

  // 2. Generar URLs firmadas temporales (1 hora)
  const { data: signed, error: signError } = await bucket.createSignedUrls(
    images.map((f) => f.name),
    60 * 60,
  );
  if (signError) throw signError;

  return images.map((file, i) => ({
    name: file.name,
    url: signed?.[i]?.signedUrl ?? '',
    size: file.metadata?.['size'],
    createdAt: file.created_at ?? undefined,
  }));
}
\end{lstlisting}

% ===========================================================================
\section{Paso 6: Rutas y navegación}
% ===========================================================================

La aplicación es de acceso libre: no hay ninguna ruta protegida, porque leer
los datos no requiere sesión. El inicio de sesión se pide solo en el momento
en que el usuario intenta modificar algo (paso 7).

\begin{lstlisting}[language=TypeScript, caption={src/app/app.routes.ts}]
// Acceso libre: cualquiera puede ver clientes y galería.
// El login se pide solo al intentar modificar datos.
export const routes: Routes = [
  {
    path: '',
    loadChildren: () => import('./tabs/tabs.routes').then((m) => m.routes),
  },
];
\end{lstlisting}

La barra de pestañas tiene dos páginas: \textbf{Clientes}
(\texttt{/tabs/tab1}) y \textbf{Galería} (\texttt{/tabs/galeria}).

\begin{lstlisting}[language=HTML, caption={src/app/tabs/tabs.page.html}]
<ion-tabs>
  <ion-tab-bar slot="bottom">
    <ion-tab-button tab="tab1" href="/tabs/tab1">
      <ion-icon aria-hidden="true" name="people"></ion-icon>
      <ion-label>Clientes</ion-label>
    </ion-tab-button>
    <ion-tab-button tab="galeria" href="/tabs/galeria">
      <ion-icon aria-hidden="true" name="images"></ion-icon>
      <ion-label>Galería</ion-label>
    </ion-tab-button>
  </ion-tab-bar>
</ion-tabs>
\end{lstlisting}

% ===========================================================================
\section{Paso 7: Inicio de sesión al modificar datos}
% ===========================================================================

Cualquier visitante puede rellenar el formulario de cliente, pero al pulsar
\textit{Guardar} (o editar o borrar un cliente) la aplicación comprueba los
permisos con el método \texttt{runAsAdmin}:

\begin{itemize}
  \item \textbf{Sin sesión}: abre una ventana de \textit{login}
        (\texttt{ion-modal}) y guarda la acción pendiente. Los datos escritos
        en el formulario no se pierden.
  \item \textbf{Con sesión de administrador}: ejecuta la acción. Si se acaba
        de iniciar sesión, se retoma la acción pendiente automáticamente.
  \item \textbf{Con sesión sin rol de administrador}: avisa de que no tiene
        permisos y no envía nada.
\end{itemize}

\begin{lstlisting}[language=TypeScript, caption={Control de permisos en src/app/tab1/tab1.page.ts}]
// Acción que el usuario quería hacer antes de que se le pidiera el login
private pendingAction: (() => Promise<void>) | null = null;

private async runAsAdmin(action: () => Promise<void>) {
  if (!this.isAuthenticated()) {
    this.pendingAction = action;
    this.loginError.set('');
    this.loginOpen.set(true);
    return;
  }
  if (!this.isAdmin()) {
    this.showToast('No tienes permisos para modificar clientes');
    return;
  }
  await action();
}

async login() {
  if (this.loginForm.invalid) {
    this.loginForm.markAllAsTouched();
    return;
  }
  try {
    const { email, password } = this.loginForm.getRawValue();
    await this.supabase.login(email, password);
    this.loginOpen.set(false);

    // Retomar lo que el usuario intentaba hacer
    const action = this.pendingAction;
    this.pendingAction = null;
    if (action) await this.runAsAdmin(action);
  } catch (e: any) {
    this.loginError.set(e?.message ?? 'No se ha podido iniciar sesión');
  }
}
\end{lstlisting}

El formulario de \textit{login} también es reactivo: el email debe tener un
formato válido y la contraseña al menos 6 caracteres. Si Supabase rechaza las
credenciales, el error se muestra dentro de la ventana.

Esta comprobación en la interfaz es solo una ayuda para el usuario. Aunque
alguien la saltara llamando directamente a la API, la política RLS
\textit{``Escritura solo para admins''} rechazaría la operación.

% ===========================================================================
\section{Paso 8: Página de gestión de clientes}
% ===========================================================================

\subsection{Lógica de la página}

La página usa \texttt{FormBuilder} para el formulario de clientes y
\textit{signals} para el estado de la vista. El mismo formulario sirve para
crear y para editar: si hay un cliente en edición
(\texttt{editingId}), se actualiza; si no, se inserta uno nuevo.

\begin{lstlisting}[language=TypeScript, caption={Fragmento de src/app/tab1/tab1.page.ts}]
readonly clients = signal<Client[]>([]);
readonly editingId = signal<string | null>(null);

readonly isAdmin = this.supabase.isAdmin;
readonly isAuthenticated = this.supabase.isAuthenticated;

readonly clientForm = this.fb.nonNullable.group({
  nombre: ['', [Validators.required, Validators.minLength(3)]],
  empresa: ['', Validators.required],
  email: ['', [Validators.required, Validators.email]],
});

async saveClient() {
  if (this.clientForm.invalid) {
    this.clientForm.markAllAsTouched();
    return;
  }
  // Solo se guarda si el usuario es admin (si no hay sesión, pide login)
  await this.runAsAdmin(async () => {
    const value = this.clientForm.getRawValue();
    const id = this.editingId();
    try {
      if (id !== null) {
        await this.supabase.updateClient(id, value);
        this.showToast('Cliente actualizado');
      } else {
        await this.supabase.addClient(value);
        this.showToast('Cliente insertado');
      }
      this.resetForm();
      await this.loadClients();
    } catch (e: any) {
      this.showToast('Error al guardar: ' + e.message);
    }
  });
}
\end{lstlisting}

\subsection{Vista con renderizado condicional}

La vista usa el \textit{control flow} de Angular. La cabecera muestra el botón
de iniciar sesión o el de cerrarla según haya sesión, y con sesión se muestra
el email y el rol. El formulario y la lista son visibles para todos; las
opciones de editar y borrar aparecen al deslizar cada cliente
(\texttt{ion-item-sliding}) y piden el \textit{login} si hace falta.

\begin{lstlisting}[language=HTML, caption={Fragmento de src/app/tab1/tab1.page.html}]
<ion-buttons slot="end">
  @if (isAuthenticated()) {
    <ion-button (click)="logout()" aria-label="Cerrar sesión">...</ion-button>
  } @else {
    <ion-button (click)="openLogin()" aria-label="Iniciar sesión">...</ion-button>
  }
</ion-buttons>
...
@if (isAuthenticated()) {
  <p>Sesión: {{ user()?.email }} · Rol: <strong>{{ userRole() ?? 'sin permisos' }}</strong></p>
}

<ion-grid>
  <ion-row>
    <!-- FORMULARIO: visible para todos, guardar exige ser admin -->
    <ion-col size="12" size-lg="5">
      @if (!isAdmin()) {
        <ion-note>Para guardar necesitas iniciar sesión como administrador.</ion-note>
      }
      <form [formGroup]="clientForm" (ngSubmit)="saveClient()">
        <ion-item>
          <ion-input label="Nombre" labelPlacement="stacked"
            formControlName="nombre" errorText="Mínimo 3 caracteres">
          </ion-input>
        </ion-item>
        <!-- ... empresa y email ... -->
      </form>
    </ion-col>

    <!-- LISTADO: lectura pública -->
    <ion-col size="12" size-lg="7">
      <ion-list>
        @for (client of clients(); track client.id) {
          <ion-item-sliding #sliding>
            <ion-item>
              <ion-label>
                <h2>{{ client.nombre }}</h2>
                <p>{{ client.empresa }} | {{ client.email }}</p>
              </ion-label>
            </ion-item>
            <ion-item-options side="end">
              <ion-item-option (click)="editClient(client, sliding)">...</ion-item-option>
              <ion-item-option color="danger" (click)="deleteClient(client, sliding)">...</ion-item-option>
            </ion-item-options>
          </ion-item-sliding>
        } @empty {
          <ion-item><ion-label>No hay clientes registrados.</ion-label></ion-item>
        }
      </ion-list>
    </ion-col>
  </ion-row>
</ion-grid>

<!-- Ventana de login: se abre al intentar modificar sin sesión -->
<ion-modal [isOpen]="loginOpen()" (didDismiss)="closeLogin()">
  <ng-template>
    <form [formGroup]="loginForm" (ngSubmit)="login()">...</form>
  </ng-template>
</ion-modal>
\end{lstlisting}

La aplicación usa siempre el tema claro: se ha desactivado la paleta oscura
automática de Ionic (\texttt{dark.system.css}) en \texttt{global.scss}.

El \texttt{ion-grid} coloca el formulario y el listado uno al lado del otro en
pantallas grandes (\texttt{size-lg}) y uno debajo del otro en el móvil
(\texttt{size="12"}).

% ===========================================================================
\section{Paso 9: Página de galería (Storage)}
% ===========================================================================

La página \texttt{src/app/galeria} muestra cada imagen del \textit{bucket}
\texttt{galeria} en una \texttt{ion-card} con la foto, el nombre del archivo,
la fecha de subida y el tamaño. Las \textit{cards} se reparten con
\texttt{ion-grid}: una columna en el móvil, dos en tablet y cuatro en
escritorio. Mientras carga se muestra un \textit{spinner} y se puede recargar
deslizando hacia abajo (\texttt{ion-refresher}).

\begin{lstlisting}[language=HTML, caption={Fragmento de src/app/galeria/galeria.page.html}]
<ion-grid>
  <ion-row>
    @for (img of images(); track img.name) {
      <!-- 1 columna en móvil, 2 en tablet, 4 en escritorio -->
      <ion-col size="12" size-sm="6" size-lg="3">
        <ion-card>
          <img [src]="img.url" [alt]="img.name" loading="lazy" />
          <ion-card-header>
            <ion-card-title>{{ img.name }}</ion-card-title>
            <ion-card-subtitle>{{ img.createdAt | date: 'dd/MM/yyyy HH:mm' }}</ion-card-subtitle>
          </ion-card-header>
          <ion-card-content>{{ img.size / 1024 | number: '1.0-1' }} KB</ion-card-content>
        </ion-card>
      </ion-col>
    } @empty {
      <ion-col size="12">No hay imágenes en el bucket "galeria".</ion-col>
    }
  </ion-row>
</ion-grid>
\end{lstlisting}

% ===========================================================================
\section{Navegación}
% ===========================================================================

La aplicación se arranca en modo desarrollo con:

\begin{lstlisting}[language=bash, caption={Arranque del servidor de desarrollo}]
ionic serve      # http://localhost:8100
\end{lstlisting}

Al entrar con \texttt{juan@gmail.com}, la cabecera muestra el email y el rol
\texttt{admin}.

\begin{figure}[H]
    \centering
    \includegraphics[width=0.75\linewidth]{image4.png}
    \caption{Usuario Juan (admin)}
    \label{fig:juan}
\end{figure}

Como administrador, Juan ve el formulario de nuevo cliente y el listado
completo, y puede editar o borrar cada cliente deslizándolo.

\begin{figure}[H]
    \centering
    \includegraphics[width=0.75\linewidth]{image5.png}
    \caption{Entrada desde admin que se puede ver la base de datos}
    \label{fig:admin}
\end{figure}

Los mismos datos aparecen en el \textit{Table Editor} de Supabase. Esto
demuestra que la aplicación lee y escribe en la base de datos real: un cliente
creado desde la aplicación aparece en Supabase, y un cambio hecho en Supabase
aparece en la aplicación al recargarla.

\begin{figure}[H]
    \centering
    \includegraphics[width=0.75\linewidth]{image6.png}
    \caption{Datos de la base de datos desde supabase}
    \label{fig:supabase-datos}
\end{figure}

Al entrar con \texttt{noeliac@gmail.com}, que no tiene perfil, la aplicación
muestra el rol \textit{sin permisos} y un aviso de acceso denegado. El bloqueo
no depende solo de la interfaz: aunque se consultara la API directamente, la
política RLS devolvería una lista vacía.

\begin{figure}[H]
    \centering
    \includegraphics[width=0.75\linewidth]{image7.png}
    \caption{Usuario que no puede ver la base de datos}
    \label{fig:sin-permisos}
\end{figure}

En la pestaña \textbf{Galería} se muestran como \textit{cards} las imágenes
subidas al \textit{bucket} \texttt{galeria} de Supabase Storage.

\begin{figure}[H]
    \centering
    \includegraphics[width=0.75\linewidth]{image8.png}
    \caption{Visualización de las imágenes en la app}
    \label{fig:galeria}
\end{figure}

% ===========================================================================
\section{Mejoras respecto al enunciado}
% ===========================================================================

Durante el desarrollo se corrigieron varios problemas del código propuesto en
el enunciado del laboratorio:

\begin{enumerate}
  \item \textbf{Conflicto de nombres.} El método CRUD se llamaba
        \texttt{createClient}, igual que la función de Supabase que crea la
        conexión. Se ha renombrado a \texttt{addClient}.
  \item \textbf{Login real.} El enunciado dejaba el \textit{login} comentado y
        suponía que había sesión. Sin sesión, RLS devuelve una lista vacía y el
        rol nunca se carga. Se ha añadido una pantalla de inicio de sesión.
  \item \textbf{Sincronización de la sesión.} Se ha añadido
        \texttt{onAuthStateChange} para que el estado se actualice al cerrar
        sesión o al refrescarse el \textit{token}.
  \item \textbf{Condición de carrera en los guards.} Los \textit{guards}
        esperan a que la sesión esté cargada antes de decidir.
  \item \textbf{Nombre engañoso.} En el enunciado \texttt{isAdmin} guardaba el
        rol (un texto), no un booleano. Ahora es un \textit{signal} calculado
        que devuelve \texttt{true} o \texttt{false}.
  \item \textbf{Operaciones bloqueadas sin aviso.} Cuando RLS bloquea un
        \texttt{update} o un \texttt{delete}, Supabase no devuelve error. Se
        añade \texttt{.select()} para detectarlo y avisar al usuario.
  \item \textbf{Componentes de Ionic actualizados.} Se ha sustituido
        \texttt{ion-label position="stacked"}, ya obsoleto, por los atributos
        \texttt{label} y \texttt{labelPlacement} de \texttt{ion-input}.
  \item \textbf{Usuarios sin permisos.} Se ha añadido un tercer nivel de
        acceso: los usuarios sin perfil no pueden ver los clientes.
\end{enumerate}

% ===========================================================================
\section{Repositorio y despliegue}
% ===========================================================================

\subsection{GitHub}

El código está publicado en \url{https://github.com/NoeCB/supabase}. Las
carpetas \texttt{node\_modules} y \texttt{www} (resultado de la compilación)
quedan fuera del repositorio gracias al archivo \texttt{.gitignore}.

\begin{lstlisting}[language=bash, caption={Subida del proyecto a GitHub}]
git add .
git commit -m "CRUD de clientes con roles (Ionic + Angular + Supabase)"
git branch -M main
git remote add origin https://github.com/NoeCB/supabase.git
git push -u origin main
\end{lstlisting}

\subsection{Vercel}

La aplicación se despliega en Vercel a partir del repositorio de GitHub. El
archivo \texttt{vercel.json} indica cómo compilarla y redirige todas las rutas
a \texttt{index.html}, para que recargar páginas como \texttt{/login} o
\texttt{/tabs/galeria} no devuelva un error 404:

\begin{lstlisting}[caption={vercel.json}]
{
  "buildCommand": "npm run build",
  "outputDirectory": "www",
  "rewrites": [{ "source": "/(.*)", "destination": "/index.html" }]
}
\end{lstlisting}

Pasos del despliegue:

\begin{enumerate}
  \item Entrar en \url{https://vercel.com} con la cuenta de GitHub.
  \item Pulsar \textit{Add New > Project} e importar el repositorio
        \texttt{NoeCB/supabase}.
  \item Dejar la configuración que lee de \texttt{vercel.json} y pulsar
        \textit{Deploy}.
  \item Al terminar, Vercel da una URL pública. Cada \texttt{git push} a la
        rama \texttt{main} vuelve a desplegar la aplicación.
\end{enumerate}

% Cambia esta URL por la que te dé Vercel
La aplicación desplegada está disponible en
\url{https://TU-PROYECTO.vercel.app}.

% ===========================================================================
\section{Conclusión}
% ===========================================================================

Se ha desarrollado una aplicación CRUD con control de acceso basado en roles y
una galería conectada a Supabase Storage. La conexión con Supabase se apoya en:

\begin{enumerate}
  \item La librería \texttt{@supabase/supabase-js} instalada en el proyecto.
  \item Las credenciales (URL y clave pública) guardadas en los
        \textit{environments}.
  \item Un servicio (\texttt{SupabaseService}) que crea el cliente y centraliza
        la autenticación, el CRUD y el acceso a Storage.
\end{enumerate}

La seguridad se aplica en dos niveles. En el \textbf{servidor}, las políticas
RLS impiden que un usuario sin rol de administrador modifique datos y que un
usuario sin perfil los lea, aunque intente hacerlo directamente contra la API.
En el \textbf{cliente}, la interfaz se adapta al rol con \texttt{@if},
ocultando las opciones que el usuario no puede usar.

% ===========================================================================
\section*{Referencias}
\addcontentsline{toc}{section}{Referencias}
% ===========================================================================

\begin{hangparas}{0.5in}{1}
Angular. (2026). \textit{Reactive forms}. \url{https://angular.dev/guide/forms/reactive-forms}

Angular. (2026). \textit{Signals}. \url{https://angular.dev/guide/signals}

Ionic. (2026). \textit{Ionic Framework documentation}. \url{https://ionicframework.com/docs}

Supabase. (2026). \textit{JavaScript client library}. \url{https://supabase.com/docs/reference/javascript}

Supabase. (2026). \textit{Row Level Security}. \url{https://supabase.com/docs/guides/database/postgres/row-level-security}

Supabase. (2026). \textit{Storage access control}. \url{https://supabase.com/docs/guides/storage/security/access-control}

Vercel. (2026). \textit{Deploying to Vercel}. \url{https://vercel.com/docs/deployments}
\end{hangparas}

\end{document}
